% Requires \usepackage{booktabs}
\begin{table}[t]
  \centering
  \caption{LIBERO-10 success rates for our reduced screening and the
  corresponding reported reference values. Higher is better.}
  \label{tab:libero10-ours-vs-reference}
  \setlength{\tabcolsep}{7pt}
  \begin{tabular}{lccc}
    \toprule
    Model & Ours (\%) & Reference (\%) & $\Delta$ (pp) \\
    \midrule
    $\pi_{0.5}$-LIBERO & 93.0 & 92.4 & +0.6 \\
    $\pi_{0.5}$-base & 0.0 & -- & -- \\
    OpenVLA-7B & 47.0 & $53.7 \pm 1.3$ & -6.7 \\
    OpenVLA-OFT & 96.0 & 94.5 & +1.5 \\
    VLA-JEPA & \textbf{98.0} & \textbf{95.8} & +2.2 \\
    SmolVLA (0.45B) & 44.0 & 71.0 & -27.0 \\
    $\pi_{0}$-FAST & 87.0 & 60.2 & +26.8 \\
    \bottomrule
  \end{tabular}

  \vspace{2pt}
  \begin{minipage}{0.96\linewidth}
    \footnotesize
    \textit{Notes.} Ours uses one seed and 10 trials for each of the 10
    LIBERO-10 tasks (100 episodes per model). Reference values use the closest
    checkpoint/input variant to the evaluated model: the released
    $\pi_{0.5}$-LIBERO checkpoint; OpenVLA-OFT with third-person and wrist
    images plus proprioception; and the 0.45B SmolVLA paper variant. The
    $\pi_{0.5}$ base model has no reported zero-shot LIBERO-10 value. Our
    SmolVLA result uses \texttt{n\_action\_steps=10}. Reported
    protocols are not uniform, so $\Delta$ is descriptive rather than a
    reproduction claim.
  \end{minipage}
\end{table}
